% S-only draft. A (Basic) is not implemented or evaluated in this fragment.
% Include in the team report, where \Owner and \MemberFour are defined.
\section{Syntactic Analysis}
\label{sec:syntactic}
\Owner{\MemberFour}{0.75}

\subsection{Idea}
S is a structural static analysis for line reachability. It asks whether the
target can be reached in the method's control-flow graph (CFG), without
evaluating branch conditions. This overapproximates executable behaviour:
every concrete execution follows CFG edges, but a CFG path may combine
contradictory conditions. S therefore returns only \texttt{unreachable} or
\texttt{unknown}, never \texttt{reachable}.

\subsection{How it works}
\paragraph{S: Structural}
Each bytecode instruction is a CFG node. Ordinary instructions have the next
instruction as successor. Conditional jumps have both the jump target and
the next instruction as successors; an unconditional \texttt{goto} has only
its target, and a return has none. Breadth-first search starts at instruction
index zero. A queue stores discovered instructions awaiting exploration, and
a visited set prevents repeated exploration. If the target is absent from
the visited set, S returns \texttt{unreachable}; otherwise it returns
\texttt{unknown}. CFG construction and search take $O(|V|+|E|)$ time and
$O(|V|+|E|)$ space when the graph is stored explicitly.

For the Group 15 dataset, we reuse the existing \texttt{javap} disassembly
loader. A small adapter maps JVM byte offsets to zero-based instruction
indices for the shared S implementation; it does not parse Java source.
The runner consumes the validator's answer-free query manifest, not expected
labels or witnesses. It checks source and class hashes, method descriptors,
and the target assignment's source-line and bytecode mapping. Missing or
ambiguous mappings and unsupported bytecode cause errors rather than verdicts.
This is a standalone \texttt{javap}-based integration, not a verified
integration with the team's concrete interpreter or \texttt{jvm2json}.

\subsection{Guarantees}
Let $R_{\mathrm{concrete}}$ contain the instructions executable for some input,
and let $R_S$ contain those reachable in the CFG. With correct target mapping
and complete control-flow edges for the supported fragment,
$R_{\mathrm{concrete}} \subseteq R_S$.
Consequently, $t \notin R_S$ proves that target $t$ is unreachable.
The converse does not hold: a graph path need not be feasible.
S is therefore sound for its unreachability claims, but incomplete.
The dataset runner accepts the restricted integer-bytecode fragment used by
the static, loop-free methods and rejects calls, exception handlers and
backward control-flow edges.

\subsection{What it cannot do}
S does not propagate variable values or solve path conditions. For example,
a path requiring both $x=-7$ and $x>-7$ remains present in the CFG, although
no input satisfies it. The actual run on all 30 queries, compiled with
JDK 17.0.21, returned 30 \texttt{unknown} and zero \texttt{unreachable}.
This covers 15 reachable and 15 unreachable ground-truth targets; S proved
none of the latter unreachable. Tests additionally exercise genuinely
disconnected target instructions using synthetic bytecode.
These results illustrate the precision limit of S on this constructed
dataset, not a failure of soundness or evidence about arbitrary Java programs.
More precise rungs must use information that S discards.
